\documentclass[12pt,a4paper]{article}

% ── Packages ──────────────────────────────────────────────────────────────────
\usepackage[margin=2.4cm]{geometry}
\usepackage{amsmath, amssymb}
\usepackage{graphicx}
\usepackage{float}
\usepackage{caption}
\usepackage[colorlinks=true, linkcolor=black!70,
            citecolor=black!60, urlcolor=black!70]{hyperref}
\usepackage{booktabs}
\usepackage{xcolor}
\usepackage{titlesec}
\usepackage{fancyhdr}
\usepackage{setspace}
\usepackage{parskip}
\usepackage{microtype}
\usepackage{enumitem}
\usepackage{mdframed}

% ── Code listing ───────────────────────────────────────────────────────────────
\usepackage{listings}
\definecolor{codebg}{RGB}{248,248,248}
\definecolor{codeframe}{RGB}{200,200,200}
\definecolor{kwcolor}{RGB}{0,100,180}
\definecolor{strcolor}{RGB}{163,21,21}
\definecolor{cmtcolor}{RGB}{100,140,100}
\definecolor{numcolor}{RGB}{180,80,0}

\lstdefinestyle{pystyle}{
  language=Python,
  backgroundcolor=\color{codebg},
  frame=single, framesep=4pt,
  rulecolor=\color{codeframe},
  basicstyle=\ttfamily\footnotesize,
  keywordstyle=\color{kwcolor}\bfseries,
  stringstyle=\color{strcolor},
  commentstyle=\color{cmtcolor}\itshape,
  numberstyle=\tiny\color{gray},
  numbers=left, numbersep=8pt,
  breaklines=true, breakatwhitespace=true,
  showstringspaces=false,
  tabsize=4,
  captionpos=b,
  xleftmargin=14pt,
  morekeywords={np,plt,RNG,abs,len,int,float,range,print,zip,append},
}
\lstset{style=pystyle}

% ── Formatting ─────────────────────────────────────────────────────────────────
\onehalfspacing
\setlength{\parskip}{5pt}

\titleformat{\section}{\large\bfseries}{\thesection}{1em}{}
  [\vspace{-3pt}\rule{\linewidth}{0.5pt}]
\titleformat{\subsection}{\normalsize\bfseries}{\thesubsection}{1em}{}

\pagestyle{fancy}
\fancyhf{}
\rhead{\small Kuramoto Model $\cdot$ Modelling Complex Systems}
\lhead{\small \leftmark}
\cfoot{\small\thepage}
\renewcommand{\headrulewidth}{0.4pt}

\mdfdefinestyle{keybox}{
  backgroundcolor=black!3,
  linecolor=black!40, linewidth=0.6pt,
  innerleftmargin=10pt, innerrightmargin=10pt,
  innertopmargin=7pt, innerbottommargin=7pt,
  roundcorner=2pt
}

% ── Title ──────────────────────────────────────────────────────────────────────
\title{
  \vspace{-0.9cm}
  {\LARGE\bfseries Synchronization Dynamics in the Kuramoto Model}\\[0.3em]
  {\large Characterization via the Complex Order Parameter}\\[0.7em]
  \rule{0.82\linewidth}{0.5pt}\\[0.4em]
  {\normalsize Modelling Complex Systems}\\[0.25em]
  {\normalsize [Student Name]\quad$\cdot$\quad[ID]\quad$\cdot$\quad[Department / University]}
}
\date{\today}

% ══════════════════════════════════════════════════════════════════════════════
\begin{document}
\maketitle
\thispagestyle{fancy}

% ── Abstract ──────────────────────────────────────────────────────────────────
\begin{mdframed}[style=keybox]
\noindent\textbf{Abstract.}\;
We study the all-to-all Kuramoto model with sine coupling as a paradigm
for collective synchronization in complex systems.
Starting from the model's equations of motion, we introduce the complex
order parameter $r(t)e^{i\psi(t)}$ as the natural measure of phase coherence
and derive the critical coupling $K_c = 2/[\pi g(0)]$ through a
self-consistent mean-field argument.
The nature of the transition — continuous, with order parameter
scaling $r \sim \sqrt{K-K_c}$ — is established analytically and confirmed
through numerical simulation with a fourth-order Runge--Kutta integrator.
Results are presented for two frequency distributions (Lorentzian and
Gaussian), and finite-size convergence is demonstrated across
$N \in \{25, 75, 150, 400\}$.
The complete simulation code is included as an appendix.
\end{mdframed}

\tableofcontents
\newpage

% ══════════════════════════════════════════════════════════════════════════════
\section{Introduction}
\label{sec:intro}

One of the most striking phenomena in complex systems is the ability of
large numbers of independently oscillating units to spontaneously
coordinate their rhythms.
This synchronization does not require a leader or external clock; it
emerges purely from the mutual interaction among oscillators.
The degree to which an individual oscillator's natural frequency differs
from its neighbours determines whether it will join the synchronized
cluster or continue to drift freely.

Yoshiki Kuramoto (1975) introduced a model that captures this competition
with remarkable economy: $N$ phase oscillators, each with its own natural
frequency, coupled to every other through a sine function.
The resulting equations are nonlinear and high-dimensional, yet they
admit an exact analytical solution in the thermodynamic limit
$N \to \infty$ via a mean-field argument — a rare property in complex
systems.

The model's significance lies in three features:
(i) the synchronization onset is a genuine second-order phase transition,
placing it in the same universality class as ferromagnetism;
(ii) it is solvable analytically, yielding an explicit expression for
the critical coupling;
(iii) it applies quantitatively to a wide range of physical and
biological systems, from power grids to neural populations.

This report derives the theory, verifies it numerically, and presents
visualizations of the full dynamics across both sides of the transition.

% ══════════════════════════════════════════════════════════════════════════════
\section{Model Formulation}
\label{sec:model}

\subsection{Equations of motion}

Let each oscillator $i = 1, \ldots, N$ have a phase $\theta_i(t) \in \mathbb{R}$
and a natural frequency $\omega_i$.
The all-to-all Kuramoto model with sine coupling reads:

\begin{mdframed}[style=keybox]
\begin{equation}
  \dot{\theta}_i = \omega_i + \frac{K}{N}
  \sum_{j=1}^{N} \sin(\theta_j - \theta_i),
  \qquad i = 1, \ldots, N.
  \label{eq:model}
\end{equation}
\end{mdframed}

The coupling constant $K \geq 0$ sets the strength of interaction.
Dividing by $N$ ensures the total force on each oscillator remains
order-1 as $N$ grows.
The sine function is the natural choice: it is $2\pi$-periodic, odd
(so the interaction between any pair is symmetric), zero when two
oscillators are perfectly in phase or in anti-phase, and maximal
at a quarter-period offset.

\subsection{Natural frequency distributions}
\label{sec:dists}

The natural frequencies $\omega_i$ are drawn i.i.d.\ from a symmetric,
unimodal distribution $g(\omega)$ centred at zero (equivalently, we work
in the co-rotating mean-frequency frame).
We study two distributions:

\paragraph{Lorentzian (Cauchy):}
\begin{equation}
  g_L(\omega) = \frac{1}{\pi}\frac{\gamma}{\omega^2 + \gamma^2},
  \qquad \gamma > 0.
  \label{eq:lorentzian}
\end{equation}
With $\gamma = 1$: peak $g_L(0) = 1/\pi$, critical coupling
$K_c^L = 2$.
The Lorentzian has heavy algebraic tails ($\sim \omega^{-2}$);
its self-consistency equation admits a closed-form solution.

\paragraph{Gaussian (Normal):}
\begin{equation}
  g_G(\omega) = \frac{1}{\sigma\sqrt{2\pi}}
  \exp\!\left(-\frac{\omega^2}{2\sigma^2}\right),
  \qquad \sigma > 0.
  \label{eq:gaussian}
\end{equation}
With $\sigma = 1$: peak $g_G(0) = 1/\sqrt{2\pi}$, critical coupling
$K_c^G = 2\sqrt{2\pi} \approx 5.01$.
The Gaussian has lighter tails; its $r(K)$ curve rises more steeply
above $K_c$ since nearly all oscillators eventually lock.

% ══════════════════════════════════════════════════════════════════════════════
\section{The Complex Order Parameter}
\label{sec:order}

\subsection{Definition}

The collective state of the oscillator population is encoded in the
\textbf{complex order parameter}:

\begin{mdframed}[style=keybox]
\begin{equation}
  r(t)\,e^{i\psi(t)} = \frac{1}{N}\sum_{j=1}^{N} e^{i\theta_j(t)}.
  \label{eq:OP}
\end{equation}
\end{mdframed}

Geometrically, each term $e^{i\theta_j}$ is a unit vector on the
complex circle at angle $\theta_j$.
Their centroid has modulus $r \in [0,1]$ (the \emph{synchronization
amplitude}) and argument $\psi$ (the \emph{mean phase}).
The two extreme cases are:
\begin{itemize}[leftmargin=2em, itemsep=2pt]
  \item $r = 0$: phases uniformly distributed; fully incoherent state.
  \item $r = 1$: all phases identical; perfectly synchronized state.
\end{itemize}

\subsection{Mean-field reduction}

Taking the imaginary part of $e^{-i\theta_i} \cdot r e^{i\psi}$ gives

\[
  r \sin(\psi - \theta_i) = \frac{1}{N}\sum_{j}\sin(\theta_j - \theta_i).
\]

Substituting back into Eq.~\eqref{eq:model}:
\begin{equation}
  \dot{\theta}_i = \omega_i + Kr\sin(\psi - \theta_i).
  \label{eq:mf}
\end{equation}

Each oscillator now couples only to the scalar mean field $(r, \psi)$,
not to $N-1$ partners individually.
Because $(r, \psi)$ is itself produced by the ensemble, Eq.~\eqref{eq:mf}
must be solved self-consistently.

% ══════════════════════════════════════════════════════════════════════════════
\section{Mean-Field Analysis of the Transition}
\label{sec:mft}

\subsection{Steady-state population structure}

In a frame rotating at the mean frequency ($\psi = 0$,
Eq.~\eqref{eq:mf} becomes $\dot\theta_i = \omega_i - Kr\sin\theta_i$).
Two populations coexist at steady state:

\begin{enumerate}[leftmargin=2.5em, itemsep=4pt]
  \item \textbf{Phase-locked} ($|\omega_i| \leq Kr$):
    a fixed point exists at $\sin\theta_i^* = \omega_i/(Kr)$,
    stable on the upper semicircle.
    These oscillators rotate at the collective frequency and
    contribute coherently to $r$.
  \item \textbf{Drifting} ($|\omega_i| > Kr$):
    no fixed point; the oscillator cycles continuously.
    By symmetry its time-averaged contribution to $\text{Re}[re^{i\psi}]$
    is zero.
\end{enumerate}

\subsection{Self-consistency and critical coupling}

The order parameter equals the contribution of the locked population.
Changing variables from $\omega$ to $\theta$ in the locked band:

\begin{equation}
  r = Kr\int_{-\pi/2}^{\pi/2}
  \cos^2\!\theta\; g(Kr\sin\theta)\,d\theta.
  \label{eq:selfcon}
\end{equation}

Non-trivial solutions ($r > 0$) first appear when the right-hand side
equals $r$ as $r \to 0^+$.
Expanding $g(Kr\sin\theta) \approx g(0)$ and integrating:

\[
  1 = Kg(0)\int_{-\pi/2}^{\pi/2}\cos^2\!\theta\,d\theta
    = \frac{\pi}{2}\,Kg(0).
\]

\begin{mdframed}[style=keybox]
\begin{equation}
  K_c = \frac{2}{\pi\,g(0)}.
  \label{eq:kc}
\end{equation}
\end{mdframed}

The critical coupling grows with the spread of natural frequencies
(lower $g(0)$) and shrinks for a more concentrated distribution
(higher $g(0)$).

\subsection{Order-parameter scaling and transition type}

Expanding Eq.~\eqref{eq:selfcon} to cubic order in $r$:
\begin{equation}
  r \;\approx\; \frac{\pi}{2}Kg(0)\,r - C(K,g)r^3 + \mathcal{O}(r^5),
\end{equation}
where $C > 0$.
Cancelling $r$ and solving gives

\begin{equation}
  r \;\sim\; \sqrt{\frac{K - K_c}{K_c}},\qquad K \gtrsim K_c,
  \label{eq:scaling}
\end{equation}

a \textbf{square-root onset} with critical exponent $\beta = 1/2$.
This is the hallmark of a \emph{continuous (second-order) phase transition}.
For the Lorentzian the self-consistency integral can be evaluated exactly,
yielding the closed-form result:

\begin{equation}
  r = \sqrt{1 - \frac{K_c}{K}},\qquad K > K_c.
  \label{eq:exact_lor}
\end{equation}

Table~\ref{tab:summary} summarises the key analytical results for both
distributions.

\begin{table}[H]
  \centering
  \caption{Summary of analytical predictions for the two frequency
           distributions used in this study.}
  \label{tab:summary}
  \renewcommand{\arraystretch}{1.35}
  \begin{tabular}{lcc}
    \toprule
    & \textbf{Lorentzian} ($\gamma=1$) & \textbf{Gaussian} ($\sigma=1$) \\
    \midrule
    $g(0)$      & $1/\pi \approx 0.318$ & $1/\sqrt{2\pi} \approx 0.399$ \\
    $K_c$       & $2.00$                & $\sqrt{2\pi} \approx 2.51$ \\[2pt]
    $r(K)$, $K > K_c$ & $\sqrt{1-K_c/K}$ (exact) & $\approx\sqrt{(K-K_c)/K_c}$ (near $K_c$) \\
    Tail decay  & Algebraic ($\omega^{-2}$) & Sub-exponential \\
    $r(K\!\to\!\infty)$ & $\to 1$ slowly & $\to 1$ faster \\
    \bottomrule
  \end{tabular}
\end{table}

% ══════════════════════════════════════════════════════════════════════════════
\section{Numerical Results}
\label{sec:results}

\subsection{Simulation method}

Equation~\eqref{eq:model} was integrated with a standard
fourth-order Runge--Kutta scheme (step $\Delta t = 0.05$).
The right-hand side coupling was computed as
$\operatorname{Im}[e^{-i\theta_i}\sum_j e^{i\theta_j}]/N$,
reducing the cost from $\mathcal{O}(N^2)$ to $\mathcal{O}(N)$ per step.
All steady-state values $\langle r \rangle$ are averages over the
final 20\% of each run; the first 80\% is discarded as transient.
The complete code is given in Appendix~\ref{app:code}.

\subsection{Phase dynamics on the unit circle}

\begin{figure}[H]
  \centering
  %% ── INSERT new_fig1_phases.png ─────────────────────────────────────
  %% 2x2 grid of unit-circle snapshots. K = 0.4, 1.8, 2.5, 6.0.
  %% Dots are HSV-coloured by phase; black arrow = order parameter.
  \includegraphics[width=0.82\textwidth]{new_fig1_phases.png}
  \caption{%
    Phase snapshots for $N=100$ oscillators (Lorentzian, $\gamma=1$,
    $K_c = 2$) at $t=50$ for four coupling values arranged in a
    $2\times2$ grid.
    Dots are coloured by their phase value using a cyclic (HSV) colourmap;
    the black arrow is the order parameter $re^{i\psi}$.
    \emph{Top-left} ($K=0.4$): uniform distribution, $r \approx 0$.
    \emph{Top-right} ($K=1.8$): slight clustering near $K_c$, small but
    nonzero $r$.
    \emph{Bottom-left} ($K=2.5$): a coherent cluster has formed above $K_c$.
    \emph{Bottom-right} ($K=6.0$): near-full synchronization; only a few
    Lorentzian-tail oscillators remain unlocked.
  }
  \label{fig:phases}
\end{figure}

\subsection{Transient and steady-state behaviour of $r(t)$}

\begin{figure}[H]
  \centering
  %% ── INSERT new_fig2_timeseries.png ────────────────────────────────
  %% r(t) for K = 0.5, 1.5, 2.0, 3.0, 5.0. Shaded transient region.
  \includegraphics[width=0.88\textwidth]{new_fig2_timeseries.png}
  \caption{%
    Time evolution of the order parameter $r(t)$ for five coupling
    strengths ($N=100$, Lorentzian).
    The grey shaded band ($t < 10$) marks the transient region; dotted
    horizontal lines show the analytical steady-state predictions from
    Eq.~\eqref{eq:exact_lor} for $K > K_c$.
    Below $K_c = 2$ the order parameter remains near zero throughout.
    Above $K_c$ it rises rapidly and saturates, with higher $K$ converging
    faster — consistent with the supercritical nature of the bifurcation.
  }
  \label{fig:timeseries}
\end{figure}

\subsection{Synchronization transition: Lorentzian and Gaussian}

\begin{figure}[H]
  \centering
  %% ── INSERT new_fig3_transition.png ───────────────────────────────
  %% Side-by-side panels: left = Lorentzian r vs K, right = Gaussian r vs K.
  %% Both N=100. Left: exact analytical curve. Right: near-Kc approx.
  \includegraphics[width=\textwidth]{new_fig3_transition.png}
  \caption{%
    Steady-state order parameter $\langle r \rangle$ as a function of
    coupling $K$ for the Lorentzian (\emph{left}) and Gaussian
    (\emph{right}) distributions ($N=100$).
    Circles are numerical data; curves are analytical predictions.
    Both distributions undergo a continuous transition at their respective
    $K_c$ values (dashed verticals), but the Lorentzian's heavier tails
    make $r$ rise more slowly and never reach 1, while the Gaussian curve
    climbs more steeply and approaches unity for large $K$.
    The shaded regions mark the incoherent (left) and synchronized (right)
    phases.
  }
  \label{fig:transition}
\end{figure}

\subsection{Finite-size effects}

\begin{figure}[H]
  \centering
  %% ── INSERT new_fig4_fsize.png ─────────────────────────────────────
  %% r vs K for N = 25, 75, 150, 400 (Lorentzian). Black dashed = theory.
  \includegraphics[width=0.82\textwidth]{new_fig4_fsize.png}
  \caption{%
    Finite-size dependence of the synchronization transition (Lorentzian,
    $\gamma=1$).
    Each coloured curve corresponds to a different system size
    $N \in \{25, 75, 150, 400\}$; the dashed black curve is the exact
    thermodynamic-limit prediction.
    For small $N$ the transition is smeared and a positive $r$ persists
    below $K_c$ due to finite-sample fluctuations of order $N^{-1/2}$.
    These artefacts decrease systematically with $N$, and the numerical
    curves converge toward the analytical prediction.
  }
  \label{fig:fsize}
\end{figure}

% ══════════════════════════════════════════════════════════════════════════════
\section{Discussion}
\label{sec:discussion}

\subsection{Anatomy of the transition}

Figures~\ref{fig:phases}--\ref{fig:transition} together paint a coherent
picture of the transition.
Below $K_c$ the system is in its incoherent fixed point: phases are
effectively random, $r$ fluctuates around $N^{-1/2}$, and oscillators
respond only to their own natural frequencies.
At $K_c$ the incoherent fixed point loses stability in a
\emph{supercritical pitchfork bifurcation}.
A new stable branch — the partially synchronized state — emerges
continuously, with the mean phase direction $\psi$ selected by
the initial conditions (spontaneous symmetry breaking of the
rotational $U(1)$ symmetry).
Above $K_c$ the locked cluster grows, which strengthens the mean field
$r$, which in turn pulls more oscillators in: a self-reinforcing cascade.

\subsection{Lorentzian vs.\ Gaussian: same physics, different scales}

Both distributions share the critical exponent $\beta = 1/2$ and the
same qualitative phase diagram.
Their differences are quantitative.
The Gaussian has $K_c \approx 2.51$ (with $\sigma=1$) versus $K_c = 2$
for the Lorentzian with $\gamma=1$.
More notably, the Gaussian's $r(K)$ curve rises more steeply above $K_c$
because its oscillators have bounded-looking frequency spreads —
far fewer oscillators reside in the extreme-frequency tails that can
never be locked.
By contrast the Lorentzian's algebraic tails always supply a drifting
population, so $r < 1$ for any finite $K$.

\subsection{Why the order parameter works}

The utility of $r e^{i\psi}$ is not merely that it compresses $N$
phases into two numbers.
Its deeper importance is that it makes the $N$-body problem self-consistent
and one-body: Eq.~\eqref{eq:mf} shows that each oscillator's dynamics
is governed solely by $r$ and $\psi$, which are themselves averages over
the ensemble.
This mean-field structure is what makes exact analysis possible and
is also what limits the model — in real sparse networks the mean field
is not uniform and the simple picture breaks down.

\subsection{Relevance to real systems}

The Kuramoto model is not merely a toy.
Power-grid generators whose rotor angles are governed by the swing
equation follow equations structurally identical to Eq.~\eqref{eq:model};
the critical coupling maps onto minimum transmission admittance required
for grid stability.
In neuroscience, the gamma-band ($\sim$40 Hz) synchrony associated with
selective attention in cortical networks is well described by Kuramoto
dynamics; the ratio $K/K_c$ can be estimated from spike-train recordings
and used to infer proximity to the critical point.

% ══════════════════════════════════════════════════════════════════════════════
\section{Conclusion}
\label{sec:conclusion}

The Kuramoto model provides one of the clearest examples of emergent
collective behaviour in a complex system: phase coherence arising
spontaneously from a disordered population of coupled oscillators.
The complex order parameter $re^{i\psi}$ is the ideal tool for
characterizing this emergence — it is zero in the incoherent phase,
grows continuously above the critical coupling $K_c = 2/[\pi g(0)]$
with universal exponent $\beta = 1/2$, and in the Lorentzian case
is given exactly by $r = \sqrt{1 - K_c/K}$.

Numerical simulations with $N = 100$ oscillators confirm these predictions
quantitatively for both Lorentzian and Gaussian frequency distributions.
Finite-size effects are mild for $N \geq 100$ and vanish as $N^{-1/2}$,
demonstrating robust convergence to the mean-field thermodynamic limit.
The side-by-side comparison of the two distributions shows that while
the value of $K_c$ and the shape of the $r(K)$ curve depend on $g$,
the second-order nature of the transition and its critical exponent
are universal.

% ══════════════════════════════════════════════════════════════════════════════
\appendix

\section{Python Simulation Code}
\label{app:code}

The following self-contained Python script produces all four figures
in this report.
It requires only \texttt{numpy} and \texttt{matplotlib}.
Run with \texttt{python kuramoto\_sim.py}; the four PNG files should be
placed in the same directory as this \LaTeX{} file before compilation.

\subsection*{Core simulation and figure generation}

\begin{lstlisting}[caption={%
  Complete Kuramoto simulation code. Generates all four figures.},
  label={lst:main}]
"""
kuramoto_sim.py  --  Kuramoto model simulation
Generates: new_fig1_phases.png, new_fig2_timeseries.png,
           new_fig3_transition.png, new_fig4_fsize.png
"""
import numpy as np
import matplotlib
matplotlib.use("Agg")
import matplotlib.pyplot as plt
from matplotlib.patches import Circle

plt.rcParams.update({
    "font.family": "serif", "font.size": 11,
    "axes.spines.top": False, "axes.spines.right": False,
})
RNG = np.random.default_rng(314)

# ── Simulation helpers ───────────────────────────────────────────────

def rhs(theta, omega, K):
    """O(N) vectorised RHS via complex exponential sum."""
    S = np.sum(np.exp(1j * theta))
    return omega + (K / len(theta)) * np.imag(np.exp(-1j * theta) * S)

def integrate(N, K, omega, T=60, dt=0.05):
    """RK4 integration. Returns (theta_history, r_array)."""
    theta = RNG.uniform(0, 2 * np.pi, N)
    steps = int(T / dt)
    r_arr   = np.empty(steps + 1)
    th_hist = np.empty((steps + 1, N))
    r_arr[0] = abs(np.mean(np.exp(1j * theta)))
    th_hist[0] = theta
    for s in range(steps):
        k1 = rhs(theta,                omega, K)
        k2 = rhs(theta + 0.5*dt*k1,   omega, K)
        k3 = rhs(theta + 0.5*dt*k2,   omega, K)
        k4 = rhs(theta +    dt*k3,     omega, K)
        theta = theta + (dt / 6) * (k1 + 2*k2 + 2*k3 + k4)
        r_arr[s+1]   = abs(np.mean(np.exp(1j * theta)))
        th_hist[s+1] = theta
    return th_hist, r_arr

def mean_r(N, K, dist="lor", T=30):
    """Steady-state order parameter (Euler, fast)."""
    omega = RNG.standard_cauchy(N) if dist=="lor" else RNG.normal(0,1,N)
    theta = RNG.uniform(0, 2*np.pi, N)
    steps = int(T / 0.05)
    acc = []
    for s in range(steps):
        S = np.sum(np.exp(1j * theta))
        theta += 0.05 * (omega + (K/N) * np.imag(np.exp(-1j*theta)*S))
        if s > int(0.8 * steps):
            acc.append(abs(np.mean(np.exp(1j * theta))))
    return float(np.mean(acc))

# ── Figure 1: 2x2 phase portraits ───────────────────────────────────

K_snap = [0.4, 1.8, 2.5, 6.0]
fig, axes = plt.subplots(2, 2, figsize=(7.5, 7.5))
axes = axes.flatten()
cmap = plt.cm.hsv

for ax, K in zip(axes, K_snap):
    omega = RNG.standard_cauchy(100)
    th_hist, r_arr = integrate(100, K, omega, T=50)
    th  = th_hist[-1];  r   = r_arr[-1]
    psi = np.angle(np.mean(np.exp(1j * th)))

    ax.set_facecolor("#FAFAFA")
    ax.add_patch(Circle((0,0), 1, fill=False,
                         color="#BBBBBB", lw=1.2, ls="--"))
    ax.axhline(0, color="#DDDDDD", lw=0.8)
    ax.axvline(0, color="#DDDDDD", lw=0.8)

    colors = cmap((th % (2*np.pi)) / (2*np.pi))
    ax.scatter(np.cos(th), np.sin(th), s=25, c=colors,
               alpha=0.85, edgecolors="white",
               linewidths=0.3, zorder=3)
    ax.annotate("", xy=(r*np.cos(psi), r*np.sin(psi)),
                xytext=(0, 0),
                arrowprops=dict(arrowstyle="-|>", color="black",
                                lw=2.0, mutation_scale=16))
    ax.set_xlim(-1.38, 1.38);  ax.set_ylim(-1.38, 1.38)
    ax.set_aspect("equal")
    ax.set_xticks([]);  ax.set_yticks([])
    for sp in ax.spines.values(): sp.set_visible(False)
    region = "K < Kc" if K < 2 else "K > Kc"
    ax.set_title(f"K={K},  r={r:.2f}  ({region})",
                 fontsize=11, pad=8)

fig.suptitle(
    "Phase snapshots  (N=100, Lorentzian, Kc=2)",
    fontsize=12, fontweight="bold")
plt.tight_layout()
plt.savefig("new_fig1_phases.png", bbox_inches="tight", dpi=150)
plt.close()

# ── Figure 2: r(t) time series ───────────────────────────────────────

fig, ax = plt.subplots(figsize=(9, 4.5))
K_ts     = [0.5, 1.5, 2.0, 3.0, 5.0]
palette  = ["#2E4057", "#A8DADC", "#FF9F1C", "#2EC4B6", "#E71D36"]
t_arr    = np.arange(1201) * 0.05

for K, col in zip(K_ts, palette):
    omega = RNG.standard_cauchy(100)
    _, r_arr = integrate(100, K, omega, T=60)
    ax.plot(t_arr, r_arr, color=col, lw=1.8, alpha=0.92,
            label=f"K={K}")
    if K > 2.0:
        ax.hlines(np.sqrt(1 - 2.0/K), 0, 60,
                  colors=col, lw=1.0, ls=":", alpha=0.6)

ax.axvspan(0, 10, alpha=0.07, color="gray")
ax.axvline(10, color="gray", lw=0.8, ls="--", alpha=0.5)
ax.text(11, 0.97, "steady state", fontsize=9, color="gray")
ax.set_xlabel("Time t"); ax.set_ylabel("r(t)")
ax.set_title("Time evolution of the order parameter",
             fontweight="bold")
ax.set_ylim(-0.03, 1.05); ax.set_xlim(0, 60)
ax.legend(loc="center right", framealpha=0.9, fontsize=10)
plt.tight_layout()
plt.savefig("new_fig2_timeseries.png", bbox_inches="tight", dpi=150)
plt.close()

# ── Figure 3: transition curves, Lorentzian and Gaussian ─────────────

Kc_lor = 2.0;  Kc_gau = np.sqrt(2 * np.pi)
K_lor  = np.arange(0, 5.5,  0.35)
K_gau  = np.arange(0, 8.5,  0.50)
r_lor  = np.array([mean_r(100, K, "lor") for K in K_lor])
r_gau  = np.array([mean_r(100, K, "gau") for K in K_gau])

K_th_l = np.linspace(Kc_lor, 5.4, 200)
r_th_l = np.sqrt(1 - Kc_lor / K_th_l)
K_th_g = np.linspace(Kc_gau, 8.4, 200)
r_th_g = np.sqrt(np.maximum(0, (K_th_g - Kc_gau) / Kc_gau))

fig, axes = plt.subplots(1, 2, figsize=(12, 4.8), sharey=True)

for ax, K_pts, r_pts, K_th, r_th, Kc, col1, col2, title in [
    (axes[0], K_lor, r_lor, K_th_l, r_th_l, Kc_lor,
     "steelblue", "firebrick",
     "Lorentzian  (gamma=1,  Kc=2.00)"),
    (axes[1], K_gau, r_gau, K_th_g, r_th_g, Kc_gau,
     "darkorange", "purple",
     f"Gaussian  (sigma=1,  Kc={Kc_gau:.2f})"),
]:
    ax.axvspan(0, Kc, alpha=0.07, color=col1)
    ax.axvspan(Kc, K_pts[-1]+1, alpha=0.07, color=col2)
    ax.scatter(K_pts, r_pts, s=45, color=col1,
               alpha=0.85, zorder=5, label="Simulation (N=100)")
    ax.plot(K_th, r_th, color=col2, lw=2.4, label="Analytical")
    ax.axvline(Kc, color="gray", lw=1.3, ls="--",
               label=f"Kc={Kc:.2f}")
    ax.set_xlabel("Coupling K"); ax.set_ylabel("<r>")
    ax.set_title(title, fontweight="bold")
    ax.set_ylim(-0.05, 1.05)
    ax.legend(fontsize=9, framealpha=0.9)

fig.suptitle(
    "Synchronization transition for two frequency distributions",
    fontsize=13, fontweight="bold")
plt.tight_layout()
plt.savefig("new_fig3_transition.png", bbox_inches="tight", dpi=150)
plt.close()

# ── Figure 4: finite-size effects ───────────────────────────────────

N_vals = [25, 75, 150, 400]; Kc = 2.0
K_pts  = np.array([0, 0.5, 1.0, 1.5, 2.0, 2.5, 3.0, 4.0, 5.0])
K_th   = np.linspace(0, 5.1, 300)
r_th   = np.where(K_th > Kc,
                  np.sqrt(1 - Kc / np.maximum(K_th, 0.01)), 0.0)
COLS   = ["#E63946", "#457B9D", "#2A9D8F", "#F4A261"]

fig, ax = plt.subplots(figsize=(8, 5))
for N, col in zip(N_vals, COLS):
    r_n = np.array([mean_r(N, K, "lor") for K in K_pts])
    ax.plot(K_pts, r_n, "o-", color=col, lw=1.8, ms=6,
            alpha=0.88, label=f"N = {N}")
ax.plot(K_th, r_th, color="black", lw=2.5, ls="--",
        label="Theory (exact)", zorder=10)
ax.axvline(Kc, color="gray", lw=1.2, ls=":", alpha=0.8)
ax.set_xlabel("Coupling strength K"); ax.set_ylabel("<r>")
ax.set_title("Finite-size convergence", fontweight="bold")
ax.set_ylim(-0.05, 1.12); ax.set_xlim(-0.1, 5.2)
ax.legend(loc="upper left", framealpha=0.9)
plt.tight_layout()
plt.savefig("new_fig4_fsize.png", bbox_inches="tight", dpi=150)
plt.close()

print("All four figures saved.")
\end{lstlisting}

\end{document}
